\documentclass[11pt,a4paper]{article}
\usepackage[utf8]{inputenc}
\usepackage[T1]{fontenc}
\usepackage[french]{babel}
\usepackage{textcomp}
\usepackage{geometry}
\geometry{margin=2.5cm}
\usepackage{fancyhdr}
\setlength{\headheight}{14.2pt}
\usepackage{tcolorbox}
\usepackage{listings}
\usepackage{xcolor}
\usepackage{setspace}
\usepackage{titlesec}
\usepackage{enumitem}
\usepackage{eurosym}
\usepackage{hyperref}
\usepackage{tikz}
\usepackage{inconsolata}
\usepackage{multicol}
\def\checkmark{\tikz\fill[scale=0.4](0,.35) -- (.25,0) -- (1,.7) -- (.25,.15) -- cycle;} 

% Entête et pied de page
\pagestyle{fancy}
\fancyhf{}
\rhead{Université de Lille - L2 Économie}
\lhead{TP Python}
\cfoot{\thepage}

% Style listings
\lstdefinelanguage{terminal}{
  morekeywords={ls,cd,pwd,mkdir,python,python3,code,dir,type,echo},
  sensitive=true
}
\lstset{
  basicstyle=\ttfamily\small,
  backgroundcolor=\color{gray!10},
  frame=single,
  language=Python,
  keywordstyle=\color{blue},
  commentstyle=\color{gray!70!black},
  stringstyle=\color{teal},
  showstringspaces=false,
  breaklines=true
}

% Interligne
\onehalfspacing

% Couleur exercices
\tcbset{
  colback=blue!5!white,
  colframe=blue!30!white,
  fonttitle=\bfseries,
  coltitle=black,
  boxrule=0.5pt,
  arc=2pt,
  left=6pt,
  right=6pt,
  top=4pt,
  bottom=4pt,
  before skip=10pt,
  after skip=10pt,
}

\title{\textbf{Premiers pas en Python: \\ Exercices complémentaires gradués}}
\author{Etienne Dagorn\thanks{\href{mailto:etienne.dagorn@univ-lille.fr}{{etienne.dagorn@univ-lille.fr}}}}
\date{}

\begin{document}
\maketitle

\tableofcontents

\newpage



\newpage
\section{Exercices Complémentaires}

\begin{tcolorbox}[title=Mode d'emploi pédagogique]
Cette banque sert à consolider les TPs, pas à remplacer la progression principale. Les niveaux indiquent le degré d'autonomie attendu.
\begin{itemize}
    \item \textbf{Niveau 1 -- Lecture} : comprendre un code existant et prédire son résultat.
    \item \textbf{Niveau 2 -- Modification} : changer légèrement un programme fourni.
    \item \textbf{Niveau 3 -- Complétion} : compléter quelques lignes manquantes.
    \item \textbf{Niveau 4 -- Construction} : écrire une solution complète à partir d'un objectif.
    \item \textbf{Niveau 5 -- Transfert} : combiner plusieurs notions dans une situation nouvelle.
\end{itemize}
Commencer par un exercice de niveau 1 ou 2 lorsqu'une notion est encore récente, puis monter progressivement.
\end{tcolorbox}

\subsection*{Lire un dictionnaire -- Niveau 1}
Avant d'écrire du code, prédire ce que ce programme affiche.
\begin{lstlisting}
prix = {"pain": 1.2, "lait": 0.95}
print(prix["pain"])
print("pates" in prix)
print(prix.get("pates", 0))
\end{lstlisting}
\begin{tcolorbox}[title=Aide]
Lire une clé existante avec \texttt{prix["pain"]}. Tester l'existence avec \texttt{in}. Utiliser \texttt{get(..., 0)} pour obtenir une valeur par défaut.
\end{tcolorbox}

\subsection*{Modifier une boucle -- Niveau 2}
Le programme affiche les valeurs positives. Le modifier pour afficher seulement les valeurs supérieures ou égales à 10.
\begin{lstlisting}
valeurs = [4, 10, 15, 2, 20]
for x in valeurs:
    if x > 0:
        print(x)
\end{lstlisting}
\begin{tcolorbox}[title=Aide]
Une seule ligne est à modifier : la condition après \texttt{if}.
\end{tcolorbox}

\subsection*{Déboguer un compteur -- Niveau 3}
Ce programme devrait compter les valeurs strictement positives, mais il ne s'exécute pas correctement. Identifier l'erreur, puis corriger.
\begin{lstlisting}
valeurs = [4, -1, 0, 8]
for x in valeurs:
    if x > 0:
        compteur = compteur + 1
print(compteur)
\end{lstlisting}
\begin{tcolorbox}[title=Aide]
Chercher quelle variable est utilisée avant d'avoir reçu une première valeur.
\end{tcolorbox}

\subsection*{Fusion de dictionnaires -- Niveau 4}
Deux paniers contiennent des quantités :
\begin{lstlisting}
p1 = {"pomme": 2, "pain": 1, "riz": 1}
p2 = {"pomme": 1, "yaourt": 4, "riz": 2}
\end{lstlisting}
\begin{itemize}
    \item Fusionne \texttt{p1} et \texttt{p2} en \texttt{p\_total} en \textbf{additionnant} les quantités lorsqu'une clé existe dans les deux.
    \item Trie \texttt{p\_total} par quantité décroissante (affichage).
\end{itemize}
\begin{tcolorbox}[title=Aide]
Crée un nouveau dict, copie \texttt{p1}, puis parcours \texttt{p2}: si la clé existe déjà, additionne.\\
Tri d'affichage : \texttt{sorted(p\_total.items(), key=lambda kv: -kv[1])}.
\end{tcolorbox}
\begin{tcolorbox}[title=Bonus -- Aide]
Dans \texttt{fusion(p,q, mode)}, gère trois cas: \texttt{"somme"} (addition), \texttt{"droite"} (remplace par \texttt{q}), \texttt{"gauche"} (garde \texttt{p}).
\end{tcolorbox}

\subsection*{Fréquence des lettres -- Niveau 4}
Demande un mot à l'utilisateur et affiche combien de fois chaque lettre apparaît (ex : \texttt{banana $\Rightarrow$ b:1, a:3, n:2}).
\begin{itemize}
    \item Affiche les lettres triées par fréquence décroissante.
\end{itemize}
\begin{tcolorbox}[title=Aide]
Parcours caractère par caractère, garde seulement les lettres (\texttt{.isalpha()}), normalise en minuscules.\\
Compteur: \texttt{d[c] = d.get(c, 0) + 1}. Tri: \texttt{sorted(d.items(), key=lambda kv: -kv[1])}.
\end{tcolorbox}
\begin{tcolorbox}[title=Bonus -- Aide]
Ignore les caractères non alphabétiques avant le comptage; filtre \texttt{a--z} via \texttt{.isalpha()} et \texttt{.lower()}.
\end{tcolorbox}

\subsection*{Classement de notes -- Niveau 3}
À partir d'un dictionnaire \texttt{\{"Alice": 14, "Bob": 10, "Clara": 16\}}, affiche :
\begin{itemize}
    \item Les noms triés par note décroissante.
    \item Le/la premier.ère et le/la dernier.ère.
\end{itemize}
\begin{tcolorbox}[title=Aide]
Trie les \texttt{items()} du dictionnaire par la \emph{valeur} décroissante: \texttt{key=lambda kv: -kv[1]}.\\
Le premier et le dernier sont accessibles par indices \texttt{[0]} et \texttt{[-1]}.
\end{tcolorbox}
\begin{tcolorbox}[title=Bonus -- Aide]
Écris \texttt{top\_k(d, k)} qui trie puis renvoie les \texttt{k} premiers éléments via un \emph{slicing} \texttt{[:k]}.
\end{tcolorbox}

% ======================================================
% BLOC 2 : EXERCICES APPLICATION / PETITS CAS RÉALISTES

\subsection*{Inventaire simple (prix \& stock) -- Niveau 5}
On représente un inventaire par :
\begin{lstlisting}
inventaire = {
  "pomme": {"prix": 0.5, "stock": 20},
  "pain":  {"prix": 1.2, "stock": 10},
  "riz":   {"prix": 2.0, "stock": 5}
}
\end{lstlisting}
\begin{itemize}
    \item Calcule la valeur totale de l'inventaire (\texttt{prix * stock}, somme).
    \item Vends une commande \texttt{\{"pomme": 3, "riz": 2\}} si le stock est suffisant ; mets à jour les stocks.
    \item Ajoute une remise de 10\% sur \texttt{"riz"} et recalcule la valeur totale.
\end{itemize}
\begin{tcolorbox}[title=Aide]
Somme totale: parcours \texttt{inventaire.items()} et ajoute \texttt{prix * stock}.\\
Avant la vente: vérifie que chaque \texttt{stock} est suffisant, puis décrémente.\\
Remise: multiplie le \texttt{prix} par \texttt{0.9}, puis recalcule la somme.
\end{tcolorbox}
\begin{tcolorbox}[title=Bonus -- Aide]
Pour trier par \emph{valeur de stock} (\texttt{prix * stock}), trie \texttt{inventaire.items()} avec \texttt{key=lambda kv: -(kv[1]["prix"]*kv[1]["stock"])}.
\end{tcolorbox}

\subsection*{Petites ``stats'' sur transactions -- Niveau 5}
On dispose d'une liste de transactions :
\begin{lstlisting}
tx = [
  {"nom":"Alice","montant":12.5},
  {"nom":"Bob","montant":7.0},
  {"nom":"Alice","montant":3.5},
  {"nom":"Clara","montant":20.0}
]
\end{lstlisting}
\begin{itemize}
    \item Calcule la dépense totale par personne (dictionnaire \texttt{nom $\rightarrow$ total}).
    \item Qui a le plus dépensé ? Affiche le classement.
\end{itemize}
\begin{tcolorbox}[title=Aide]
Cumule dans un dict: \texttt{totaux[nom] = totaux.get(nom,0) + montant}.\\
Classement: \texttt{sorted(totaux.items(), key=lambda kv: -kv[1])}.
\end{tcolorbox}
\begin{tcolorbox}[title=Bonus -- Aide]
Stocke pour chaque nom \texttt{nb} d'achats et \texttt{total}; la moyenne est \texttt{total/nb}.
\end{tcolorbox}

\subsection*{Fonctions utilitaires (liste de nombres) -- Niveau 4}
Écris des fonctions :
\begin{itemize}
    \item \texttt{moyenne(L)}, \texttt{mediane(L)} (liste triée/tri à faire), \texttt{mode(L)} (valeur la plus fréquente).
    \item Teste sur plusieurs listes (\texttt{[1,2,2,3]}, \texttt{[10]}, \texttt{[]}).
\end{itemize}
\begin{tcolorbox}[title=Aide]
\textbf{Moyenne} : \texttt{sum/len} avec test liste vide.\\
\textbf{Médiane} : trie une \emph{copie}; pair $\Rightarrow$ moyenne des deux centraux.\\
\textbf{Mode} : dict de fréquences, \texttt{max()} sur les valeurs; si ex aequo, renvoyer une liste.
\end{tcolorbox}
\begin{tcolorbox}[title=Bonus -- Aide]
Pense à tester les cas limites: liste vide, un seul élément, multimodal.
\end{tcolorbox}

\subsection*{Groupements simples (sans pandas) -- Niveau 5}
On dispose d'un catalogue \texttt{[{catégorie, produit, prix}]} (liste de dicts).
\begin{itemize}
    \item Calcule, par \textbf{catégorie}, le \textbf{nombre} de produits et le \textbf{prix moyen}.
    \item Affiche les catégories triées par prix moyen décroissant.
\end{itemize}
\begin{tcolorbox}[title=Aide]
Crée un dict \texttt{groupes[cat] = \{nb, somme\}} en parcourant le catalogue.\\
Puis calcule \texttt{prix\_moy = somme/nb} et regroupe dans une liste de tuples à trier.
\end{tcolorbox}

\subsection*{Nettoyage léger -- Niveau 3}
À partir d'une liste de prénoms \texttt{[" alice ", "Bob", "ALICE", "bob ", "  clara"]} :
\begin{itemize}
    \item Supprime espaces superflus, harmonise en \texttt{Nom propre} (majuscule initiale seulement).
    \item Produit une liste \texttt{unique} triée.
\end{itemize}
\begin{tcolorbox}[title=Aide]
Écris une petite fonction de nettoyage: \texttt{strip()}, \texttt{lower()}, \texttt{capitalize()}.\\
Construis la liste \emph{unique} en testant \texttt{if x not in ...} puis trie avec \texttt{sorted()}.
\end{tcolorbox}

\newpage

\subsection*{Sauvegarder les résultats dans un fichier (100\% Python) -- Niveau 5}

Dans cet exercice, vous allez apprendre à \textbf{enregistrer automatiquement les résultats de vos exercices précédents dans un fichier texte} en utilisant uniquement Python, sans passer par le terminal.

\begin{tcolorbox}[title=Avant de commencer]
Vérifiez que vous savez déjà créer un fichier \texttt{.py}, l'enregistrer, l'exécuter et retrouver le fichier créé dans le dossier courant. Cet exercice introduit \texttt{open()}, qui est plus avancé que les sorties avec \texttt{print()}.
\end{tcolorbox}

\begin{itemize}
    \item Créez un nouveau script nommé \texttt{tp2\_solutions.py}.
    \item En haut du script, importez le module \texttt{datetime} et ouvrez un fichier texte pour y écrire les résultats :
\begin{lstlisting}[language=Python]
from datetime import datetime

# Creation (ou ecrasement) du fichier de resultats
fichier_resultats = open("resultats_tp2.txt", "w", encoding="utf-8")

# Fonction utilitaire : ecrit a l'ecran et dans le fichier
def log(texte):
    print(texte)
    fichier_resultats.write(str(texte) + "\n")

log("=== Resume des exercices du TP2 ===")
log("Execute le : " + str(datetime.now()))
log("")  # ligne vide
\end{lstlisting}

    \item Pour chaque exercice du TP2 (1 à 15), ajoutez au moins un exemple dont le résultat sera enregistré dans le fichier en utilisant \texttt{log()} :
\begin{lstlisting}[language=Python]
# Exemple : Exercice 2
notes = [12, 14, 10, 8]
moyenne = sum(notes) / len(notes)
log("Ex2 - Moyenne des notes : " + str(round(moyenne, 1)))
\end{lstlisting}

    \item Une fois tous les résultats enregistrés, n'oubliez pas de fermer le fichier pour le sauvegarder :
\begin{lstlisting}[language=Python]
fichier_resultats.close()
\end{lstlisting}

\end{itemize}

\begin{tcolorbox}[title=À retenir]
Après exécution de \texttt{tp2\_solutions.py}, un fichier \texttt{resultats\_tp2.txt} sera créé automatiquement dans le même dossier.  
Il contiendra toutes les sorties principales de vos exercices.  
Vous n'avez pas besoin d'utiliser le terminal : tout se fait directement en Python.
\end{tcolorbox}



\end{document}
